% ============================================================
%  cap00_capa.tex — Capa e Ficha Técnica
%  Guia de Fontes de Dados para Avaliação de Políticas Públicas
% ============================================================

\begin{titlepage}
\newgeometry{top=3cm, bottom=2.5cm, left=3cm, right=2.5cm}
\thispagestyle{empty}
\color{clearblue}

% Série
\vspace*{1cm}
{\centering
{\large\bfseries\color{clearlightblue}
\MakeUppercase{Série: Avaliação na Prática}}\par}

\vspace{1.8cm}

% Título principal
{\centering
{\fontsize{36}{42}\selectfont\bfseries\color{clearblue}
Guia de Fontes de Dados}\par}

\vspace{0.4cm}

{\centering
{\fontsize{24}{30}\selectfont\bfseries\color{clearblue}
para Avaliação de}\par}

\vspace{0.2cm}

{\centering
{\fontsize{24}{30}\selectfont\bfseries\color{clearblue}
Políticas Públicas}\par}

\vspace{2cm}

% Instituições
{\centering
{\large\color{clearblue}
FGV CLEAR \,|\, FGV EESP \,|\, Global Evaluation Initiative}\par}

\vspace{0.5cm}

% Data
{\centering{\normalsize\color{cleargray}[Mês/Ano]}\par}

\vspace{2cm}

% Linha separadora
{\color{clearlightblue}\rule{\linewidth}{0.8pt}}

\vspace{0.6cm}

% Resumo/sinopse
\begin{tcolorbox}[
  enhanced, colback=clearlightbluelight!50, colframe=clearblue,
  sharp corners, boxrule=0.5pt, left=8pt, right=8pt, top=6pt, bottom=6pt
]
\small\color{black}
Este guia documenta o ecossistema de fontes de dados disponíveis para a avaliação de
políticas públicas no Brasil. Cobre pesquisas amostrais, registros administrativos,
censos, bases internacionais e ferramentas emergentes (APIs, \textit{web scraping} e
inteligência artificial), organizados por área temática: mercado de trabalho, saúde,
educação, assistência social, segurança pública, infraestrutura, finanças públicas e
fontes internacionais. Para cada fonte são apresentados produtor, cobertura, periodicidade,
unidade de análise, pontos fortes, limitações e indicações de uso em avaliação de impacto.
\end{tcolorbox}

\restoregeometry
\end{titlepage}

% ── Ficha Técnica ──────────────────────────────────────────
\thispagestyle{empty}
\vspace*{2cm}

{\large\bfseries\color{clearblue} Ficha Técnica}

\bigskip

\noindent\textbf{Organizadores:} [Nome 1], [Nome 2]

\smallskip
\noindent\textbf{Equipe técnica:} [Listar membros]

\bigskip

\noindent\textbf{Apresentação.} Este guia integra a série de publicações
\textit{Avaliação na Prática}, desenvolvida pelo FGV CLEAR com o objetivo de
ampliar o acesso a conhecimentos sobre monitoramento e avaliação com foco em
políticas públicas.

\medskip

\noindent Fundado em 2015, o FGV CLEAR tem se dedicado ao fortalecimento da cultura
de gestão orientada por evidências no Brasil e em países lusófonos. Com sede na
Escola de Economia de São Paulo da Fundação Getulio Vargas (FGV EESP), o FGV CLEAR
atua como centro regional da Iniciativa CLEAR
(\textit{Centers for Learning on Evaluation and Results}).

\medskip

\noindent Saiba mais sobre o FGV CLEAR e acesse outras publicações em:
\textbf{www.fgvclear.org}

\vfill

\noindent{\small\color{cleargray}%
Os textos desta publicação são de responsabilidade exclusiva dos autores.
As opiniões expressas não refletem necessariamente a posição das instituições parceiras.}

\clearpage
